The example of zero-shot baseline for the question "What is AI agent?":
\begin{quote}
        \label{appendix:zero_shot_without_context_answer}
\textbf{Understanding AI Agents}

\textbf{Prerequisites for Understanding "AI Agent"}

To understand what an AI Agent is, let's first look at the prerequisites—the building blocks that help explain the topic better. Think of prerequisites like the LEGO bricks you need before building a cool spaceship!

\textbf{1. What is "Artificial Intelligence" (AI)?}

\begin{itemize}
    \item \textbf{Simple Idea:} AI is about making computers smart—enabling them to act, think, or learn like humans do.
    \item \textbf{What it means:} It's not just a calculator. AI tries to solve problems, recognize objects in pictures, or understand language. (Think Siri or Alexa—that's basic AI!)
\end{itemize}

\subsubsection*{2. What is a "Program" or "Software"?}

\begin{itemize}
    \item \textbf{Simple Idea:} A program is a set of step-by-step instructions that a computer follows to perform a task.
    \item \textbf{What it means:} If you tell a computer "Add 2 and 3," that's a simple instruction. Most software consists of very long, complex instructions.
\end{itemize}

\subsubsection*{3. What is a "Goal" or "Objective"?}

\begin{itemize}
    \item \textbf{Simple Idea:} Every task needs a finish line. If you go for a walk, your goal might be to reach the park.
    \item \textbf{What it means:} In technology, an AI Agent needs a specific goal to achieve (e.g., "Book me a flight to London" or "Keep the temperature at 72 degrees").
\end{itemize}

\subsection*{Now, Let's Define "AI Agent"}

Now that we have the building blocks (AI, Instructions, Goals), we can put them together!

\begin{quote}
\textbf{An AI Agent is basically a smart computer program that can \textit{sense} its environment, \textit{think} about what to do based on its goal, and then \textit{take action} to get closer to that goal, all by itself.}
\end{quote}

\subsubsection*{Analogy: A Smart Robot Cleaner}

Imagine a robot vacuum cleaner:

\begin{enumerate}
    \item \textbf{Sense (Perceive):} It has sensors that "see" the dirt, walls, and furniture in the room. (This is the input from the environment.)
    \item \textbf{Think (Process):} It runs an AI program that thinks: "Goal is to clean the whole floor. I see dirt here, so I should go there." (This is the "brain" deciding what to do.)
    \item \textbf{Act (Execute):} It moves, sucks the dirt, and changes direction. (This is the output/action.)
\end{enumerate}

The key difference between a simple program and an AI Agent is \textbf{autonomy} (self-control).

\begin{itemize}
    \item \textbf{Simple Program:} You have to tell it \textit{every single move}. ("Move forward 1 foot. Turn 90 degrees. Suck dirt.")
    \item \textbf{AI Agent:} You just tell it the \textit{goal}. ("Clean the whole room.") It figures out the steps on its own!
\end{itemize}

\subsection*{Summary}

\begin{tabular}{|p{3cm}|p{10cm}|}
\hline
\textbf{Feature} & \textbf{What it Means in Plain English} \\
\hline
\textbf{Agent} & Something that acts or operates on its own. \\
\hline
\textbf{AI} & The "brain" that allows it to learn, decide, and understand things. \\
\hline
\textbf{Combined (AI Agent)} & A smart, self-directing program that takes in information, decides on a plan to reach a goal, and executes that plan. \\
\hline
\end{tabular}

\subsubsection*{Real-World Examples:}

\begin{itemize}
    \item A chatbot that doesn't just answer questions, but \textit{books} your meeting.
    \item A self-driving car (it senses the road, decides to brake, and executes the braking action).
    \item An AI trading bot that monitors the stock market and automatically buys or sells when it thinks the price is right.
\end{itemize}
\end{quote}
